\documentclass[11pt]{article}
\usepackage[margin=0.9in]{geometry}
\usepackage{amsmath,amssymb,amsthm,booktabs,mathrsfs}
\usepackage[T1]{fontenc}
\usepackage{lmodern,microtype}
\usepackage[colorlinks=true,linkcolor=blue,urlcolor=blue]{hyperref}
\newtheorem{theorem}{Theorem}[section]
\newtheorem{lemma}[theorem]{Lemma}
\newtheorem{proposition}[theorem]{Proposition}
\newtheorem{corollary}[theorem]{Corollary}
\theoremstyle{remark}
\newtheorem{remark}[theorem]{Remark}
\newcommand{\norm}[1]{\left\|#1\right\|}
\title{Toward a four-dimensional Yang--Mills construction\\
\large Notebook 20: the first Wilson small-field integration}
\author{Research working document}
\date{17 October 2026}
\begin{document}
\maketitle
\begin{abstract}
The physical source-coordinate construction of Notebook 19 and
the overlap expansion of Notebook 18 yield an actual first-step
small-field integration of the non-Abelian Wilson action.
For the bare step, every action term is local before and after
the exact pivot substitution. This verifies the support-valued
precision and interaction hypotheses directly. The retained
sharp cutoff is fixed; its nonlinear Jacobian, Haar density,
gauge remainder and normalization terms are kept explicitly.
The normalized fluctuation logarithm has connected activities
whose anchored norm is bounded by
$C[gR^3+\exp(-cR^2)]$, for a logarithmic cutoff $R$.
Constants are uniform in periodic volume at fixed blocking
ratio. Strict support here is in the independent current
background. Its transport to the next coarse field, the
subsequent quantum iterations, and the complementary
large-field regions remain the next construction steps.
\end{abstract}

This notebook applies the internally checked finite-regulator
results in Notebooks 15--19. It is a proof of a specified
small-field contribution. Completing the full existence
argument requires that contribution to be incorporated into
an iteration that controls the whole regulator measure and
the continuum observable hierarchy.

